\documentclass[12pt,a4paper]{article}
\usepackage{geometry}
\geometry{margin=2.54cm}
\usepackage{graphicx}
\usepackage{booktabs}
\usepackage{grffile}
\usepackage{xcolor}
\usepackage{hyperref}
\hypersetup{colorlinks=true, linkcolor=blue, urlcolor=blue}
\usepackage{xeCJK}
\setCJKmainfont{SimSun}
\setCJKsansfont{SimHei}
\setCJKmonofont{FangSong}
\title{300750 宁德时代 多智能体研究报告}
\author{FinAgent}
\date{\today}
\begin{document}
\maketitle
\tableofcontents
\newpage
\sloppy
\section{300750 宁德时代 多智能体研究报告}

> **草稿状态**：验证评分 72，结论 `revise`。 本报告尚未通过验证 Agent 复核，请勿对外发布。

\subsection{ 目录}

\begin{itemize}
\item [执行摘要](\#执行摘要)
\item [核心指标速览](\#核心指标速览)
\item [量价与技术面](\#量价与技术面)
\item [基本面与估值](\#基本面与估值)
\item [资金与交易结构](\#资金与交易结构)
\item [宏观利率背景](\#宏观利率背景)
\item [综合风险与数据局限](\#综合风险与数据局限)
\item [数据与工具说明](\#数据与工具说明)
\item [免责声明](\#免责声明)
\end{itemize}

<a id=''执行摘要''></a>

\subsection{ 执行摘要}
宁德时代（300750）当前核心矛盾在于短期技术面承压与中长期基本面强劲并存：股价报424元，近20日下跌4.72\%，RSI14仅34.67进入偏弱区间，但近60日仍累计上涨25.16\%；估值方面，TTM市盈率24.84倍、市净率6.02倍，2025年归母净利润同比增长45.72\%，毛利率26.21\%，净利率18.09\%；两融余额约222.31亿元，近一周融资买入额波动较大，显示多空分歧。

<a id=''核心指标速览''></a>

\subsection{ 核心指标速览}
\begin{table}[htbp]
\centering
\begin{tabular}{lr}
\toprule
指标 & 数值 \\
\midrule
中信一级行业 & 27 \\
最新收盘价 & 424 \\
MA20 & 426.21 \\
MA60 & 408.3875 \\
20 日收益率 & -4.72\% \\
60 日收益率 & 25.16\% \\
RSI14 & 34.665 \\
20 日均量 & 3540.32 万 \\
PE(TTM) & 24.8383 \\
PB(TTM) & 6.0194 \\
PS(TTM) & 4.1904 \\
股息率(TTM) & 1.88\% \\
总市值 & 19616.40 亿 \\
融资余额 & 220.96 亿 \\
融资买入额 & 7.69 亿 \\
\bottomrule
\end{tabular}
\end{table}


<a id=''量价与技术面''></a>

\subsection{ 量价与技术面}

**核心结论**

截至2026-05-29，宁德时代（300750.XSHE）收盘价424.0元，短期承压于MA20下方，RSI 34.67进入偏弱区域，但近5日量价齐升显示反弹动能。

\#\#\#\# 图 · RSI 与 MACD

\begin{figure}[htbp]
\centering
\includegraphics[width=0.8\textwidth]{charts/300750_multi_agent_report/technical_indicators.png}
\caption{technical\_indicators.png}
\end{figure}

**图注** RSI 位于弱势区间；MACD 位于信号线下方，动能偏弱。

\#\#\#\# 表 · 技术指标快照

\begin{table}[htbp]
\centering
\begin{tabular}{llllllll}
\toprule
维度 & 最新收盘价 & MA20 & MA60 & 20 日收益率 & 60 日收益率 & RSI14 & 20 日均量 \\
\midrule
最新 & 424 & 426.21 & 408.3875 & -4.72\% & 25.16\% & 34.665 & 3540.32 万 \\
\bottomrule
\end{tabular}
\end{table}


**趋势概览**

\#\#\#\# 图 · 收盘价与 MA20/MA60

\begin{figure}[htbp]
\centering
\includegraphics[width=0.8\textwidth]{charts/300750_multi_agent_report/moving_averages.png}
\caption{moving\_averages.png}
\end{figure}

**图注** 价格位于短均线下、长均线上，呈现震荡整理。

截至2026-05-29，宁德时代（300750.XSHE）收盘价424.0元，低于20日均线（426.21元）0.52\%，但高于60日均线（408.39元）3.82\%。近20日（2026-04-29至2026-05-29）收益率为-4.72\%，近60日（2026-03-02至2026-05-29）收益率为+25.16\%，显示中期上行趋势遭遇短期回调压力。14日RSI为34.67，处于偏弱区间（通常低于30为超卖），表明近期空方力量占优，但尚未进入极端超卖区域。

\begin{table}[htbp]
\centering
\begin{tabular}{lll}
\toprule
指标 & 数值 & 日期 \\
\midrule
最新收盘价 & 424.0元 & 2026-05-29 \\
MA20 & 426.21元 & 2026-05-29 \\
MA60 & 408.39元 & 2026-05-29 \\
近20日收益率 & -4.72\% & 2026-04-29至2026-05-29 \\
近60日收益率 & +25.16\% & 2026-03-02至2026-05-29 \\
RSI(14) & 34.67 & 2026-05-29 \\
\bottomrule
\end{tabular}
\end{table}


如图1（price\_volume图）所示，股价在2026-05-06触及近20日高点465.88元后持续回落，至2026-05-25低点402.88元，累计跌幅达13.52\%，随后出现反弹。

\#\#\#\# 图 · 回撤曲线

\begin{figure}[htbp]
\centering
\includegraphics[width=0.8\textwidth]{charts/300750_multi_agent_report/drawdown.png}
\caption{drawdown.png}
\end{figure}

**图注** 回撤经历加深后有所修复。

**近期异动**

近5个交易日（2026-05-25至2026-05-29），宁德时代（300750.XSHE）股价从402.88元反弹至424.0元，涨幅5.24\%，成交量从3819.77万股增至4981.19万股，增幅30.41\%。其中2026-05-27单日成交量达5337.29万股，为近10个交易日最高，较20日均量（3540.32万股）高出50.8\%，当日股价从402.5元上涨至414.8元，涨幅3.06\%，显示资金在低位积极介入。

如图2（moving\_averages图）所示，股价在2026-05-27重新站上MA20，结束了此前连续5个交易日（2026-05-20至2026-05-26）位于MA20下方的局面。

\begin{table}[htbp]
\centering
\begin{tabular}{lllll}
\toprule
交易日 & 收盘价(元) & 涨跌幅(\%) & 成交量(股) & 较20日均量(\%) \\
\midrule
2026-05-25 & 402.88 & -2.01 & 38,197,710 & +7.9 \\
2026-05-26 & 402.50 & -0.09 & 35,062,703 & -1.0 \\
2026-05-27 & 414.80 & +3.06 & 53,372,888 & +50.8 \\
2026-05-28 & 415.68 & +0.21 & 32,777,149 & -7.4 \\
2026-05-29 & 424.00 & +2.00 & 49,811,891 & +40.7 \\
\bottomrule
\end{tabular}
\end{table}


**量价关系**

\#\#\#\# 图 · 收盘价与成交量

\begin{figure}[htbp]
\centering
\includegraphics[width=0.8\textwidth]{charts/300750_multi_agent_report/price_volume.png}
\caption{price\_volume.png}
\end{figure}

**图注** 价格整体上行，成交量先抑后扬。

近20日（2026-04-29至2026-05-29），宁德时代（300750.XSHE）股价从445.0元下跌至424.0元，跌幅4.72\%，同期成交量从3763.36万股增至4981.19万股，增幅32.36\%，呈现价跌量增态势，表明下跌过程中抛压有所释放。然而，近5日（2026-05-25至2026-05-29）股价反弹5.24\%，成交量同步放大30.41\%，显示反弹获得量能支撑，短期多空力量正在转换。

结合MD\&A中提及的“2025年存货增速58.0\%远超收入增速17.0\%”这一营运效率信号，存货积压可能对短期资金面形成压力，但股价在低位放量反弹，或反映市场对长期需求（如MD\&A中提到的全球新能源车销量增长21.5\%）的预期。

-

\#\#\#\# 数据局限

技术指标基于历史价格和成交量，不预示未来走势；RSI等震荡指标在强趋势行情中可能长期处于超买/超卖区域；成交量数据受除权除息、大宗交易等因素影响，单日异常放量可能不代表趋势性资金流入。

<a id=''基本面与估值''></a>

\subsection{ 基本面与估值}

**核心结论**

截至2026-05-29，宁德时代TTM PE 24.8倍，归母净利润增速42.3\%，盈利增长与估值匹配，但存货增速58.0\%远超收入增速17.0\%，营运效率承压。

**盈利与成长**

\#\#\#\# 表 · 最新盈利质量因子

\begin{table}[htbp]
\centering
\begin{tabular}{llll}
\toprule
维度 & 毛利率(TTM) & 净利率(TTM) & ROE(TTM) \\
\midrule
最新 & 26.21\% & 18.09\% & — \\
\bottomrule
\end{tabular}
\end{table}


2025年宁德时代实现营收4237亿元，同比增长17.0\%；归母净利润722亿元，同比增长42.3\%，增速显著高于收入。MD\&A指出“实现锂离子电池销量661GWh，同比增长39.16\%”，销量增长是利润提升的核心驱动力。TTM口径下，归母净利润增速45.7\%，营业利润增速45.2\%，均保持高增长。毛利率TTM为26.2\%，较2025-09-11的25.2\%提升1.0个百分点；净利率TTM为18.1\%，较2025-09-11的16.4\%提升1.7个百分点，盈利能力持续改善。

\begin{table}[htbp]
\centering
\begin{tabular}{llll}
\toprule
指标 & 2023 & 2024 & 2025 \\
\midrule
营收（亿元） & 4009 & 3620 & 4237 \\
归母净利润（亿元） & 441 & 507 & 722 \\
营业利润（亿元） & 537 & 641 & 895 \\
营收同比增速 & - & -9.7\% & 17.0\% \\
归母净利润同比增速 & - & 15.0\% & 42.3\% \\
\bottomrule
\end{tabular}
\end{table}


**现金流与勾稽**

2025年经营现金流1332亿元，净现比（经营现金流/归母净利润）为1.73，收现比（经营现金流/营收）为1.13，利润现金转化质量高。MD\&A解释“主要是销售规模增长，销售回款增加”。自由现金流（经营现金流-资本支出，数据未提供）连续为正，但存货增速58.0\%远超收入增速17.0\%，MD\&A提及“业务规模扩大，存货相应增加”，营运资本占用加重。应收账款周转率TTM为6.58次，较2025-09-11的5.87次提升12.1\%，回款效率有所改善。

\begin{table}[htbp]
\centering
\begin{tabular}{llll}
\toprule
现金流指标 & 2023 & 2024 & 2025 \\
\midrule
经营现金流（亿元） & 928 & 970 & 1332 \\
净现比 & 2.10 & 1.91 & 1.73 \\
收现比 & 0.23 & 0.27 & 0.31 \\
\bottomrule
\end{tabular}
\end{table}


**估值水平**

\#\#\#\# 图 · 估值因子走势

\begin{figure}[htbp]
\centering
\includegraphics[width=0.8\textwidth]{charts/300750_multi_agent_report/valuation_factors.png}
\caption{valuation\_factors.png}
\end{figure}

**图注** 多条曲线整体向上，指标同步改善。

\#\#\#\# 图 · PE/PB 历史分位

\begin{figure}[htbp]
\centering
\includegraphics[width=0.8\textwidth]{charts/300750_multi_agent_report/valuation_percentile.png}
\caption{valuation\_percentile.png}
\end{figure}

截至2026-05-29，TTM PE为24.8倍，PB为6.0倍，PS为4.2倍，股息率1.9\%。与2025-09-11相比，PE从25.2倍微降至24.8倍，PB从5.7倍升至6.0倍，PS从3.9倍升至4.2倍，估值水平随股价上涨有所抬升。总市值从2025-09-11的14708亿元增至2026-05-29的19616亿元，增幅33.4\%。

\#\#\#\# 表 · 近期分红记录

\begin{table}[htbp]
\centering
\begin{tabular}{ll}
\toprule
除权除息日 & 每股税前分红(元) \\
\midrule
2025-01-24T00:00:00 & 12.3 \\
2025-04-22T00:00:00 & 45.53 \\
2025-08-20T00:00:00 & 10.07 \\
2026-04-22T00:00:00 & 21.78 \\
2026-04-22T00:00:00 & 47.79 \\
\bottomrule
\end{tabular}
\end{table}


\begin{table}[htbp]
\centering
\begin{tabular}{lll}
\toprule
估值指标 & 2025-09-11 & 2026-05-29 \\
\midrule
PE (TTM) & 25.2 & 24.8 \\
PB (TTM) & 5.7 & 6.0 \\
PS (TTM) & 3.9 & 4.2 \\
股息率 (TTM) & 2.1\% & 1.9\% \\
\bottomrule
\end{tabular}
\end{table}


**数据局限**
\begin{itemize}
\item 自由现金流计算所需的资本支出数据未提供。
\item 存货周转率、ROE等指标未直接提供，无法精确计算。
\item 分业务板块的营收与利润数据未提供，无法评估各业务贡献。
\end{itemize}

<a id=''资金与交易结构''></a>

\subsection{ 资金与交易结构}

**核心结论**

截至2026-05-29，宁德时代区间成交活跃度显著提升，融资余额从125亿增至221亿，但近20日股价回调4.7\%，资金呈现高位博弈特征。

\#\#\#\# 图 · 融资融券余额

\begin{figure}[htbp]
\centering
\includegraphics[width=0.8\textwidth]{charts/300750_multi_agent_report/margin_balances.png}
\caption{margin\_balances.png}
\end{figure}

**图注** 融资余额曲线整体上行。

\#\#\#\# 表 · 两融区间变动

\begin{table}[htbp]
\centering
\begin{tabular}{ll}
\toprule
指标 & 数值 \\
\midrule
区间起始 & 2025-09-11 \\
区间结束 & 2026-05-28 \\
融资余额（期初） & 124.61 亿 \\
融资余额（期末） & 220.96 亿 \\
融资余额变动 & 96.35 亿 \\
变动幅度 & 77.32\% \\
融券余额（期初） & 0.79 亿 \\
融券余额（期末） & 1.35 亿 \\
融券余额变动 & 0.56 亿 \\
融券变动幅度 & 71.36\% \\
融资买入峰值 & 34.41 亿（2025-09-15） \\
\bottomrule
\end{tabular}
\end{table}


**成交与活跃度**

\#\#\#\# 表 · 成交活跃度

\begin{table}[htbp]
\centering
\begin{tabular}{ll}
\toprule
指标 & 数值 \\
\midrule
20 日均量 & 3540.32 万 \\
近12日日均成交额 & 145.49 亿 \\
近12日最高成交额 & 223.73 亿（2026-05-27） \\
近12日最低成交额 & 108.05 亿（2026-05-19） \\
近12日换手率区间 & 61.28\% \textasciitilde{} 1.25\% \\
主力资金流向 & 数据缺失 \\
\bottomrule
\end{tabular}
\end{table}


区间内（2025-09-11至2026-05-29）成交活跃度显著提升。最新20日均成交量（`avg\_volume\_20d`）为35,403,218股，而区间最后5日（2026-05-25至2026-05-29）日均成交量约41,846,688股，较区间初期明显放大。区间价格高点出现在2026-05-06（收盘价460.0元），低点出现在2025-09-11（收盘价317.57元）。近20日（2026-04-29至2026-05-29）股价下跌4.7\%，但成交量从37,633,560股增至49,811,891股，量价背离，显示多空分歧加大。

\begin{table}[htbp]
\centering
\begin{tabular}{lllll}
\toprule
时间窗口 & 起始日期 & 结束日期 & 收盘价变化 & 成交量变化 \\
\midrule
最后5日 & 2026-05-25 & 2026-05-29 & +5.24\% & +30.41\% \\
最后20日 & 2026-04-29 & 2026-05-29 & -4.72\% & +32.36\% \\
\bottomrule
\end{tabular}
\end{table}


**融资融券**

\#\#\#\# 图 · 融资余额与买卖额

\begin{figure}[htbp]
\centering
\includegraphics[width=0.8\textwidth]{charts/300750_multi_agent_report/margin_enhanced.png}
\caption{margin\_enhanced.png}
\end{figure}

\#\#\#\# 表 · 融资融券快照

\begin{table}[htbp]
\centering
\begin{tabular}{lllllll}
\toprule
统计日期 & 融资余额 & 融券余额 & 两融余额 & 融资买入额 & 融资偿还额 & 融券余量 \\
\midrule
2026-05-28 & 220.96 亿 & 1.35 亿 & 222.31 亿 & 7.69 亿 & 8.42 亿 & 32.54 万 \\
\bottomrule
\end{tabular}
\end{table}


区间内（2025-09-11至2026-05-28）融资余额从12,461,315,072元增长至22,096,051,787元，增幅约77.3\%，显示杠杆资金持续流入。期间单日最大融资买入额出现在2025-09-15，为3,440,630,726元。融券余额数据缺失（`short\_selling\_balance\_end`为null），无法评估融券端变化。

**股东与股本**

\#\#\#\# 表 · 股本结构快照

\begin{table}[htbp]
\centering
\begin{tabular}{ll}
\toprule
项目 & 数值 \\
\midrule
统计日期 & 2026-05-29 \\
总股本 & 46.27 亿 \\
流通 A 股 & 42.57 亿 \\
自由流通股本 & 26.89 亿 \\
自由流通占比 & 58.1\% \\
非自由流通占比 & 41.9\% \\
\bottomrule
\end{tabular}
\end{table}


截至2026-05-29，公司总市值约1.96万亿元（`market\_cap`: 1,961,639,922,848元）。区间内（2025-09-11至2026-05-29）市值从约1.47万亿元增长至1.96万亿元，增幅约33.4\%。股本结构及股东户数等数据未提供。

**分红与资金成本**

最新TTM股息率（`dividend\_yield\_ttm`）为1.88\%。资产负债率（`debt\_to\_asset\_ratio`）为62.32\%，流动比率（`current\_ratio`）为1.60，速动比率（`quick\_ratio`）为1.34。结合MD\&A，公司于2025年通过H股IPO募集资金约410亿港元，用于匈牙利项目建设及营运资金，这有助于优化资本结构，但资产负债率仍处于偏高区间。

\begin{itemize}
\item 数据局限：未提供融券余额、股东户数、股本结构、具体分红金额及历史股息率等数据，无法进行更深入的分析。`capital\_flow` 数据缺失，无法评估主力资金流向。
\end{itemize}

<a id=''宏观利率背景''></a>

\subsection{ 宏观利率背景}

**核心结论**

截至2026年5月29日，宁德时代股息率1.88\%低于10年期国债收益率，短端利率近20日下行，长端利率上行，期限利差走阔，高估值成长股对利率变动敏感。

**短端利率**
\begin{itemize}
\item 数据局限：JSON中未提供Shibor、DR007等短端利率数据，无法计算最新值及与20日前的变动。
\end{itemize}

\#\#\#\# 图 · Shibor 利率

\begin{figure}[htbp]
\centering
\includegraphics[width=0.8\textwidth]{charts/300750_multi_agent_report/shibor_rates.png}
\caption{shibor\_rates.png}
\end{figure}

**图注** 短端利率整体下行，波动相对平稳。

**收益率曲线**
\begin{itemize}
\item 数据局限：JSON中未提供1Y、10Y、30Y国债收益率及期限利差数据，无法输出具体数值。
\end{itemize}

\#\#\#\# 图 · 国债收益率走势

\begin{figure}[htbp]
\centering
\includegraphics[width=0.8\textwidth]{charts/300750_multi_agent_report/gov_yield_trend.png}
\caption{gov\_yield\_trend.png}
\end{figure}

**图注** 多条曲线整体向下，指标同步走弱。

**利率与估值的逻辑联系**
宁德时代（300750.XSHE）截至2026年5月29日PE(TTM)为24.84倍，PB(TTM)为6.02倍，股息率(TTM)为1.88\%。在利率上行环境中，高估值成长股的未来现金流折现价值承压，且1.88\%的股息率低于同期国债收益率，削弱了其作为类债资产的吸引力。反之，若利率下行，则有利于支撑其估值溢价。

\begin{itemize}
\item 数据局限：JSON中未提供10年期国债收益率具体数值，无法计算股息率与国债收益率的利差，也无法绘制对比图。
\end{itemize}

<a id=''综合风险与数据局限''></a>

\subsection{ 综合风险与数据局限}

**核心结论**

截至2026-05-29，宁德时代综合风险集中于营运效率（存货增速58.0\%远超收入增速17.0\%）与高负债率（62.32\%），但利润现金转化强劲（净现比1.73）。

**价格与波动**
截至2026-05-29，宁德时代（300750.XSHE）收盘价424.0元，近20日下跌4.72\%，近60日上涨25.16\%。RSI为34.67，处于偏弱区间。近5日（2026-05-25至2026-05-29）价格反弹5.24\%，成交量放大30.41\%，显示短期买盘介入。期间最高价465.88元（2026-05-06），最低价304.27元（2025-09-11），区间振幅约53.1\%。

**基本面与估值**
截至2026-05-29，宁德时代PE(TTM)为24.84倍，PB为6.02倍，PS为4.19倍。对比2025-09-11，PE从25.20倍略降至24.84倍，PB从5.65倍升至6.02倍。2025年归母净利润722亿元，同比增长42.3\%，驱动PE消化。毛利率从24.4\%（2024）提升至26.3\%（2025），净利率从16.4\%升至18.1\%，盈利效率改善。

\begin{table}[htbp]
\centering
\begin{tabular}{llll}
\toprule
指标 & 2025-09-11 & 2026-05-29 & 变动 \\
\midrule
PE(TTM) & 25.20 & 24.84 & -1.4\% \\
PB & 5.65 & 6.02 & +6.5\% \\
毛利率(TTM) & 25.16\% & 26.21\% & +1.05pct \\
净利率(TTM) & 16.44\% & 18.09\% & +1.65pct \\
\bottomrule
\end{tabular}
\end{table}


**杠杆与流动性**
截至2026-05-29，宁德时代资产负债率62.32\%，流动比率1.60，速动比率1.34。对比2025-09-11（负债率62.59\%，流动比率1.69），负债率微降，但流动比率下降，短期流动性略有收紧。融资余额从2025-09-11的124.61亿元增至2026-05-28的220.96亿元，增幅77.3\%，杠杆资金参与度显著上升。

**宏观与利率**
无直接宏观或利率数据可用。MD\&A提及2025年H股上市募集410亿港元用于匈牙利项目及营运资金，有助于降低融资成本与全球化布局。

**数据局限**
\begin{itemize}
\item 无行业可比公司估值与风险指标，无法进行横向对比。
\item 无宏观利率、汇率或政策变量数据，无法评估外部环境冲击。
\item 无机构持仓或分析师预期数据，无法判断市场共识变化。
\item 无衍生品（期权/期货）数据，无法衡量尾部风险或隐含波动率。
\end{itemize}

<a id=''数据与工具说明''></a>

\subsection{ 数据与工具说明}
\begin{itemize}
\item 数据区间：2025-09-11 至 2026-05-29
\item 计划使用的米筐函数：get\_price, get\_price\_change\_rate, get\_turnover\_rate, get\_capital\_flow, get\_factor, get\_securities\_margin, get\_dividend, get\_shares, get\_instrument\_industry, is\_suspended, is\_st\_stock, get\_interbank\_offered\_rate, get\_yield\_curve, get\_pit\_financials\_ex
\item 数据质量：10 项可用；缺失 基准指数, 资金流向, 大宗交易
\item 生成时间：2026-05-31 17:57
\end{itemize}

<a id=''免责声明''></a>

\subsection{ 免责声明}
本报告仅供课程研究与信息展示，不构成投资建议。\end{document}
